\documentclass[11pt]{article}
\usepackage[a4paper,margin=1in]{geometry}
\usepackage{amsmath,amssymb,mathtools}
\usepackage{booktabs,array}
\usepackage{enumitem}
\usepackage[dvipsnames]{xcolor}
\usepackage{hyperref}
\usepackage{microtype}
\usepackage{comment}

\hypersetup{colorlinks=true,linkcolor=MidnightBlue,urlcolor=MidnightBlue,
  pdftitle={Recitation 6: Correlation, Independence, and Causation},
  pdfauthor={Xuanxi Zhang}}

% Default: student handout. Define \WITHSOLUTIONS before inputting this file
% to include full answers and the instructor's 75-minute pacing guide.
\newif\ifsolutions
\ifdefined\WITHSOLUTIONS
  \solutionstrue
\else
  \solutionsfalse
\fi

\newcommand{\Pbb}{\mathbb{P}}
\newcommand{\E}{\mathbb{E}}
\newcommand{\Var}{\operatorname{Var}}
\newcommand{\Cov}{\operatorname{Cov}}
\newcommand{\rv}[1]{\widetilde{#1}}
\newcommand{\indep}{\mathrel{\perp\!\!\!\perp}}
\newcommand{\lmmse}{\ell_{\mathrm{MMSE}}}

\ifsolutions
  \newenvironment{answer}{%
    \par\medskip\noindent\textcolor{MidnightBlue}{\textbf{Solution.}}\ }
    {\par\medskip}
  \newcommand{\workingspace}[1]{}
\else
  \excludecomment{answer}
  \newcommand{\workingspace}[1]{\vspace{#1}}
\fi

\newcommand{\keybox}[1]{%
  \begin{center}
  \fcolorbox{MidnightBlue}{MidnightBlue!5}{%
    \begin{minipage}{0.91\textwidth}#1\end{minipage}}
  \end{center}}
\setlist[enumerate]{label=(\alph*),leftmargin=2em,itemsep=0.5em}
\setlist[itemize]{leftmargin=1.5em,itemsep=0.25em}

\title{\huge Recitation 6\\[0.35em]
  \normalsize Correlation, Independence, and Causation\\[0.35em]
  \large\ifsolutions Solutions and Teaching Notes\else Review and Practice Problems\fi}
\author{Xuanxi Zhang}
\date{October 16, 2026}

\begin{document}
\maketitle
\noindent\textit{Reading:}
\href{https://github.com/cfgranda/ps4ds/blob/main/ps4ds_preprint.pdf}
{\emph{Probability and Statistics for Data Science}},
Sections 8.2--8.6 and 8.8.

\section{Review: covariance and correlation}
Assume finite second moments. Write
$\mu_{\rv{x}}=\E[\rv{x}]$ and $\sigma_{\rv{x}}^2=\Var(\rv{x})$.
\[
  \Cov(\rv{x},\rv{y})
  =\E[(\rv{x}-\mu_{\rv{x}})(\rv{y}-\mu_{\rv{y}})]
  =\E[\rv{x}\rv{y}]-\mu_{\rv{x}}\mu_{\rv{y}}.
\]
If both variances are positive, the \textbf{correlation coefficient} is
\[
  \boxed{\rho_{\rv{x},\rv{y}}
    =\frac{\Cov(\rv{x},\rv{y})}{\sigma_{\rv{x}}\sigma_{\rv{y}}}.}
\]
Correlation measures the direction and strength of a \emph{linear}
association. It is dimensionless; covariance depends on the units.

\paragraph{Properties to remember.}
\begin{itemize}
  \item \textbf{Symmetry and bounds:}
    $\rho_{\rv{x},\rv{y}}=\rho_{\rv{y},\rv{x}}$ and
    $-1\leq\rho_{\rv{x},\rv{y}}\leq1$.
  \item \textbf{Shifts and rescaling:} for nonzero constants $b,d$,
    \[
      \Cov(b\rv{x}+a,d\rv{y}+c)=bd\,\Cov(\rv{x},\rv{y}),\qquad
      \rho_{b\rv{x}+a,d\rv{y}+c}
      =\operatorname{sign}(bd)\rho_{\rv{x},\rv{y}}.
    \]
    Positive changes of scale leave correlation unchanged; reversing one
    variable's sign reverses the correlation.
  \item \textbf{Perfect correlation:} $|\rho|=1$ exactly when
    $\rv{y}=\alpha+\beta\rv{x}$ with probability one, for some
    $\beta\ne0$. The sign of $\rho$ is the sign of $\beta$.
  \item \textbf{Variance of a sum:}
    $\Var(\rv{x}+\rv{y})=\Var(\rv{x})+\Var(\rv{y})
      +2\Cov(\rv{x},\rv{y})$.
\end{itemize}

\paragraph{Connection to linear prediction.}
The best \emph{affine} predictor of $\rv{y}$ from $\rv{x}=x$ is
\[
  \lmmse(x)=\mu_{\rv{y}}
     +\frac{\Cov(\rv{x},\rv{y})}{\Var(\rv{x})}(x-\mu_{\rv{x}})
  =\mu_{\rv{y}}+\rho_{\rv{x},\rv{y}}
     \frac{\sigma_{\rv{y}}}{\sigma_{\rv{x}}}(x-\mu_{\rv{x}}).
\]
Its MSE is $(1-\rho_{\rv{x},\rv{y}}^2)\sigma_{\rv{y}}^2$.
Thus $\rho^2$ is the fraction of variance explained by the best affine
predictor. If $\rho=0$, that predictor is just the mean; nonlinear
prediction may still help. A constant variable has zero covariance with
every variable, but its correlation coefficient is \emph{undefined}.

\newpage
\section{Mini-lecture: uncorrelated versus independent}
\textbf{Uncorrelated} means $\Cov(\rv{x},\rv{y})=0$.
\textbf{Independent} means that, for every pair of events defined by the
two variables,
\[
  \Pbb(\rv{x}\in A,\rv{y}\in B)
  =\Pbb(\rv{x}\in A)\Pbb(\rv{y}\in B).
\]
For discrete variables this is equivalent to
$p_{\rv{x},\rv{y}}(x,y)=p_{\rv{x}}(x)p_{\rv{y}}(y)$ for every $x,y$.

\paragraph{Independence implies uncorrelation.}
With finite second moments, independence gives
\[
  \E[\rv{x}\rv{y}]=\E[\rv{x}]\E[\rv{y}]
  \quad\Longrightarrow\quad \Cov(\rv{x},\rv{y})=0.
\]
The reverse implication is false: covariance checks only one mixed moment,
whereas independence constrains the entire joint distribution.

\paragraph{Worked example: a nonlinear relationship with zero correlation.}
Let $\rv{x}$ take $-1,0,1$ with equal probability, and set
$\rv{y}=\rv{x}^{\,2}$. The possible pairs are
\[
\begin{array}{c|ccc}
 (x,y)&(-1,1)&(0,0)&(1,1)\\ \hline
 \Pbb(\rv{x}=x,\rv{y}=y)&1/3&1/3&1/3
\end{array}
\]
By symmetry,
\[
  \E[\rv{x}]=0,\qquad \E[\rv{y}]=\frac23,\qquad
  \E[\rv{x}\rv{y}]=\frac{-1+0+1}{3}=0.
\]
Hence $\Cov(\rv{x},\rv{y})=0$, and $\rho_{\rv{x},\rv{y}}=0$
because both variables have positive variance. However,
\[
  \Pbb(\rv{x}=0,\rv{y}=0)=\frac13
  \ne\frac19=\Pbb(\rv{x}=0)\Pbb(\rv{y}=0).
\]
They are dependent: observing $\rv{x}$ determines $\rv{y}$ exactly.
The positive and negative contributions to covariance cancel.

\keybox{\textbf{The Gaussian exception.}
If $(\rv{x},\rv{y})$ is \emph{jointly Gaussian}, then
uncorrelated $\Longleftrightarrow$ independent.
It is not enough that each variable separately has a Gaussian distribution.}

\section{Review: correlation versus causation}
Association describes how variables vary together. A causal question asks
what would happen if we \emph{intervened} to change one variable.
For example, rain can cause both umbrella use and wet sidewalks:
\[
  \text{Umbrella use}\ \longleftarrow\ \text{Rain}\
  \longrightarrow\ \text{Wet sidewalks}.
\]
Their association does not establish that using an umbrella makes a
sidewalk wet. A common cause is a \textbf{confounder}; reverse causation
can also produce association. Randomized assignment can support a causal
comparison. Observational adjustment requires additional causal assumptions,
as in Recitation 5. Neither a large $|\rho|$ nor a small $|\rho|$ alone
determines a causal effect; ``variance explained'' is a predictive statement.

\newpage
\section*{Practice problems}
\noindent Three core problems; Problem 4 is a backup for a fast class.
The indicated times are for independent work before discussion.

\subsection*{Problem 1: Covariance properties (Exercise 8.2)}
\noindent\textit{Independent work: 10 minutes.}
Let $\rv{a}$ and $\rv{b}$ have finite variances. Prove the following:
\begin{enumerate}
  \item For any constants $\alpha,\beta\in\mathbb{R}$,
    \[
      \Cov(\beta\rv{a}+\alpha,\rv{b})
      =\beta\Cov(\rv{a},\rv{b}).
    \]
    Conclude that if $\rv{a}$ and $\rv{b}$ are uncorrelated, so are
    $\beta\rv{a}+\alpha$ and $\rv{b}$.
  \item The covariance satisfies
    \[
      -\sqrt{\Var(\rv{a})\Var(\rv{b})}
      \leq\Cov(\rv{a},\rv{b})
      \leq\sqrt{\Var(\rv{a})\Var(\rv{b})}.
    \]
    Include the case where one variance is zero.
  \item The sum and difference satisfy
    \[
      \Cov(\rv{a}+\rv{b},\rv{a}-\rv{b})
      =\Var(\rv{a})-\Var(\rv{b}).
    \]
\end{enumerate}
\workingspace{9cm}

\begin{answer}
Define centered random variables
$\rv{a}_0=\rv{a}-\E[\rv{a}]$ and
$\rv{b}_0=\rv{b}-\E[\rv{b}]$.

\textbf{(a)} Since
$\beta\rv{a}+\alpha-\E[\beta\rv{a}+\alpha]=\beta\rv{a}_0$,
\[
  \Cov(\beta\rv{a}+\alpha,\rv{b})
  =\E[\beta\rv{a}_0\rv{b}_0]
  =\beta\Cov(\rv{a},\rv{b}).
\]
If the original covariance is zero, the transformed covariance is zero,
including when $\beta=0$.

\textbf{(b)} Cauchy--Schwarz gives
\[
  |\Cov(\rv{a},\rv{b})|
  =|\E[\rv{a}_0\rv{b}_0]|
  \leq\sqrt{\E[\rv{a}_0^2]\E[\rv{b}_0^2]}
  =\sqrt{\Var(\rv{a})\Var(\rv{b})}.
\]
This proves both bounds. If a variance is zero, its centered variable is
zero with probability one, so the covariance is also zero. The bound still
holds, although correlation is then undefined.

\textbf{Alternative proof of (b).} If $\Var(\rv{b})>0$, set
$t=\Cov(\rv{a},\rv{b})/\Var(\rv{b})$. Then
\[
  0\leq\E[(\rv{a}_0-t\rv{b}_0)^2]
    =\Var(\rv{a})-\frac{\Cov(\rv{a},\rv{b})^2}{\Var(\rv{b})}.
\]
Rearrange and take square roots; handle zero variance as above.

\textbf{(c)} Center the sum and difference, then multiply:
\begin{align*}
  \Cov(\rv{a}+\rv{b},\rv{a}-\rv{b})
    &=\E[(\rv{a}_0+\rv{b}_0)(\rv{a}_0-\rv{b}_0)]\\
    &=\E[\rv{a}_0^2-\rv{b}_0^2]
     =\boxed{\Var(\rv{a})-\Var(\rv{b})}.
\end{align*}
No independence assumption is needed. Equal variances make the sum and
difference uncorrelated; independence does not follow in general.
\end{answer}

\newpage
\subsection*{Problem 2: An affine function (Exercise 8.11)}
\noindent\textit{Independent work: 7 minutes.}
The random variables $\rv{a}$ and $\rv{b}$ satisfy
\[
  \mu_{\rv{a}}=2,\qquad \mu_{\rv{b}}=-2,\qquad
  \sigma_{\rv{a}}^2=4,\qquad \sigma_{\rv{b}}^2=9,\qquad
  \rho_{\rv{a},\rv{b}}=0.25.
\]
Define $\rv{c}=2\rv{b}+3$. Find the linear minimum mean-squared-error
predictor of $\rv{c}$ given $\rv{a}=a$.

\medskip
\noindent\textit{Suggested steps:} Convert correlation to covariance;
find $\E[\rv{c}]$ and $\Cov(\rv{a},\rv{c})$; then identify the slope
and intercept of the predictor.
\workingspace{12cm}

\begin{answer}
The standard deviations are $\sigma_{\rv{a}}=2$ and
$\sigma_{\rv{b}}=3$, so
\[
  \Cov(\rv{a},\rv{b})
  =\rho_{\rv{a},\rv{b}}\sigma_{\rv{a}}\sigma_{\rv{b}}
  =\frac14\cdot2\cdot3=\frac32.
\]
Linearity of expectation and the covariance rule from Problem 1 give
\[
  \mu_{\rv{c}}=2(-2)+3=-1,\qquad
  \Cov(\rv{a},\rv{c})
  =2\Cov(\rv{a},\rv{b})=3.
\]
Consequently,
\[
  \boxed{\lmmse(a)
  =\mu_{\rv{c}}+\frac{\Cov(\rv{a},\rv{c})}{\Var(\rv{a})}
      (a-\mu_{\rv{a}})
  =-1+\frac34(a-2)=\frac34a-\frac52.}
\]

\paragraph{Check by transforming the predictor of $\rv{b}$.}
First predict $\rv{b}$ by
$-2+\frac38(a-2)$. Multiplying this prediction by $2$ and adding $3$
gives the same answer. The predictor passes through the point
$(\mu_{\rv{a}},\mu_{\rv{c}})=(2,-1)$.

\paragraph{Optional discussion: how much does the prediction improve?}
Here $\Var(\rv{c})=2^2\cdot9=36$, and the positive rescaling of
$\rv{b}$ leaves the correlation unchanged:
$\rho_{\rv{a},\rv{c}}=1/4$. Therefore
\[
  \mathrm{MSE}_{\mathrm{linear}}
  =\left(1-\frac1{16}\right)36
  =\frac{135}{4}=33.75.
\]
The constant predictor $-1$ has MSE $36$, so the affine predictor explains
$1/16=6.25\%$ of the variance. This is optional; the original exercise
only asks for the predictor.

\paragraph{Teaching point.}
Use standard deviations to recover covariance from correlation, but use
the input variance $\Var(\rv{a})$ in the slope denominator. The answer
minimizes MSE over affine functions. It need not equal
$\E[\rv{c}\mid\rv{a}=a]$; the exercise does not assume a jointly
Gaussian distribution or an affine conditional mean.
\end{answer}

\newpage
\subsection*{Problem 3: Ice cream sales (Exercise 8.15)}
\noindent\textit{Independent work: 8 minutes.}
Alma is analyzing data for an ice cream company. The marketing team
argues that a strong correlation between advertising expenses and sales
shows that its advertising is effective. Alma obtains the following
\emph{sample covariance matrix}:
\begin{center}
\begin{tabular}{@{}lrrr@{}}
\toprule
 & Advertising & Sales & Temperature\\
\midrule
Advertising & 900 & 960 & 540\\
Sales & 960 & 1600 & 720\\
Temperature & 540 & 720 & 400\\
\bottomrule
\end{tabular}
\end{center}
\begin{enumerate}
  \item Compute the sample correlation between advertising expenses and sales.
  \item Should Alma accept the team's conclusion? Use the temperature
    entries to inform your explanation, and distinguish association from
    a causal effect.
\end{enumerate}
\noindent\textit{Reminder:} Divide a sample covariance by the product of
the sample standard deviations. The diagonal entries above are variances.
\workingspace{11cm}

\begin{answer}
\textbf{(a)} The sample standard deviations of advertising, sales, and
temperature are $\sqrt{900}=30$, $\sqrt{1600}=40$, and
$\sqrt{400}=20$. Thus
\[
  \boxed{\text{Sample correlation of advertising and sales}
    =\frac{960}{30\cdot40}=0.8.}
\]
This is a strong positive linear association in the observed data.

\textbf{(b)} Both quantities are even more strongly associated with
temperature:
\begin{align*}
  \text{Advertising--temperature correlation}
    &=\frac{540}{30\cdot20}=0.9,\\
  \text{Sales--temperature correlation}
    &=\frac{720}{40\cdot20}=0.9.
\end{align*}
A plausible explanation is that hotter weather increases demand for ice
cream, and the company also spends more on advertising in hotter weather.
Temperature could therefore be a common cause of advertising and sales:
\[
  \text{Advertising}\ \longleftarrow\ \text{Temperature}\
  \longrightarrow\ \text{Sales}.
\]
This is a possible causal explanation, not something established by the
covariance matrix itself.

Alma can agree that advertising and sales are positively associated, but
should \textbf{not accept the claim that the association proves the ads
increase sales}. The matrix alone does not determine whether the causal
effect of advertising is positive, zero, or negative.

\paragraph{Discussion: what evidence would help?}
Randomly assign different advertising levels to comparable markets or
time periods, so assignment is independent of the potential sales outcomes,
and compare sales. Avoid spillovers between experimental groups.
With observational data, adjusting for temperature may help, but a causal
interpretation additionally requires assumptions such as no remaining
confounding and adequate overlap.

\paragraph{Common mistakes.}
The denominator is $30\cdot40$, not $900\cdot1600$.
The value $0.8$ does not mean that ads cause $80\%$ of sales.
Its square, $0.64$, describes the sample variance explained by a simple
OLS regression with an intercept, not the fraction of sales caused by ads.
Also, identifying a plausible confounder does not prove that ads have no effect.
\end{answer}

\newpage
\subsection*{Problem 4 (backup): A point in a disk}
\noindent\textit{Independent work: 5 minutes; use only if the class is ahead.}
Choose a point uniformly from the unit disk. Its coordinates have joint pdf
\[
  f_{\rv{x},\rv{y}}(x,y)=
  \begin{cases}
    1/\pi,& x^2+y^2\leq1,\\
    0,& \text{otherwise}.
  \end{cases}
\]
\begin{enumerate}
  \item Without evaluating integrals, use symmetry to find
    $\E[\rv{x}]$, $\E[\rv{y}]$, and $\E[\rv{x}\rv{y}]$.
    Are the coordinates uncorrelated?
  \item Are the coordinates independent? Justify your answer using the
    disk's support or a pair of events.
\end{enumerate}
\workingspace{11cm}

\begin{answer}
\textbf{(a)} Reflecting the disk across the vertical axis preserves the
distribution and changes the sign of $\rv{x}$, so $\E[\rv{x}]=0$.
Reflection across the horizontal axis similarly gives $\E[\rv{y}]=0$.
Reflecting across either axis reverses the sign of the product, hence
$\E[\rv{x}\rv{y}]=0$. Therefore
\[
  \boxed{\Cov(\rv{x},\rv{y})=0,\qquad\rho_{\rv{x},\rv{y}}=0.}
\]
Both variances are finite and positive because the coordinates are bounded
and nonconstant.

\textbf{(b)} Let $A=\{\rv{x}>3/4\}$ and $B=\{\rv{y}>3/4\}$.
Each event has positive probability: each corresponds to a region of
positive area inside the disk. But if both occurred, then
\[
  \rv{x}^{\,2}+\rv{y}^{\,2}
     >\frac9{16}+\frac9{16}=\frac98>1,
\]
which is impossible. Thus
\[
  \Pbb(A\cap B)=0\ne\Pbb(A)\Pbb(B).
\]
The coordinates are \textbf{dependent but uncorrelated}. Geometrically,
once one coordinate is close to the edge of the disk, the other must be
close to zero.
\end{answer}

\vfill
\noindent\textit{Source:} adapted from the 2025 archive,
\texttt{Recitations/Recitation 1 Problems.pdf}, Problem 2 [3.31].
The covariance question has been added.

\ifsolutions
\newpage
\section*{Instructor pacing guide (75 minutes)}
\begin{center}
\begin{tabular}{@{}lll@{}}
\toprule
Time & Activity & Suggested use\\
\midrule
0--10 min & Correlation review & Definition, properties, affine prediction\\
10--20 min & Uncorrelation vs. independence & Teach the worked counterexample\\
20--25 min & Correlation vs. causation & Common causes; intervention\\
25--42 min & Exercise 8.2 & 10 min individual work, 7 min discussion\\
42--55 min & Exercise 8.11 & 7 min individual work, 6 min discussion\\
55--71 min & Exercise 8.15 & 8 min individual work, 8 min discussion\\
71--75 min & Wrap-up & Revisit the two invalid implications\\
\bottomrule
\end{tabular}
\end{center}
If the class finishes early, use Problem 4 with five minutes of independent
work followed by discussion. Keep it optional; do not shorten the three
core work periods to fit it in.

\paragraph{Board plan for the mini-lecture.}
Write ``independent $\Rightarrow$ uncorrelated'' and prove it using
factorization of the product expectation. Then ask whether the reverse
holds. Work through the three-point example on page 2: list the pairs,
compute the covariance, and disprove independence with one event pair.
Emphasize that $\rv{y}=\rv{x}^{\,2}$ is perfectly predictable even though
its best affine predictor is constant. For a numerical check,
$\Var(\rv{x})=2/3$ and $\Var(\rv{y})=2/9$.

\paragraph{Why the Gaussian exception needs ``jointly.''}
For a jointly Gaussian pair, zero covariance makes the joint density
factor into the product of its marginals (when both variances are positive).
If needed, explain why Gaussian marginals alone are insufficient: take
$\rv{z}\sim\mathcal N(0,1)$ and an independent fair sign
$\rv{s}\in\{-1,1\}$. Then $\rv{z}$ and $\rv{s}\rv{z}$ are each standard
Gaussian and have covariance zero, but their absolute values are identical.
They are dependent and are not jointly Gaussian.

\paragraph{A short derivation of the affine predictor, if needed.}
For a predictor $\alpha+\beta\rv{x}$, centering gives
\[
\begin{split}
  \E[(\rv{y}-\alpha-\beta\rv{x})^2]
  ={}&(\mu_{\rv{y}}-\alpha-\beta\mu_{\rv{x}})^2
      +\sigma_{\rv{y}}^2\\
     &-2\beta\Cov(\rv{x},\rv{y})+\beta^2\sigma_{\rv{x}}^2.
\end{split}
\]
First set $\alpha=\mu_{\rv{y}}-\beta\mu_{\rv{x}}$; then minimize the
quadratic in $\beta$, giving $\beta=\Cov(\rv{x},\rv{y})/\sigma_{\rv{x}}^2$.
Substitution gives the MSE formula on page 1.

\paragraph{Closing check.}
Ask students to explain why each statement can fail:
``Zero correlation means independence'' and ``High correlation proves a
causal effect.'' A nonlinear relationship addresses the first; a common
cause addresses the second. Ask what extra assumption makes the first true
and what study design can help with the second.

\paragraph{Sources and adaptations.}
The review follows Carlos Fernandez-Granda's
\emph{Probability and Statistics for Data Science}, Sections 8.2--8.6 and
8.8. Problems 1--3 rephrase Exercises 8.2, 8.11, and 8.15, retaining all
their numerical data and requested conclusions; scaffolding has been added.
The backup reuses the disk-independence question from the 2025 archive
and adds a covariance calculation by symmetry. Its source handout is
dated January 31, 2025.
\fi

\end{document}
